% Eight graders, one protocol. One claim: whatever grades your rollouts, the
% same instruments attach. Eight distinct grader archetypes each present the
% same small set of verbs, so an instrument binds to the interface, not the
% grader. The rlSignal accent marks the grader family; a single bus collects
% all eight into one protocol panel on the right.
\documentclass[border=8pt]{standalone}
% House preamble for every reward-lens TikZ figure.
% The figure file sets \documentclass{standalone} and, before it is \input, the build
% sets \rltheme (light|dark). This file loads the matching palette and defines the shared
% visual vocabulary, so consistency across figures comes from here, not from discipline.
%
% Engine: lualatex or xelatex gives the real Inter (matches the site). pdflatex falls back
% to Fira Sans, which is close. Either way the figures do not speak Computer Modern.

\usepackage{iftex}
\ifLuaTeX
  \usepackage{fontspec}\setsansfont{Inter}[BoldFont={Inter SemiBold}, ItalicFont={Inter Italic}]
\else\ifXeTeX
  \usepackage{fontspec}\setsansfont{Inter}[BoldFont={Inter SemiBold}, ItalicFont={Inter Italic}]
\else
  \usepackage[T1]{fontenc}\usepackage[sfdefault]{FiraSans}
\fi\fi
\renewcommand{\familydefault}{\sfdefault}

\usepackage{amsmath}
\usepackage{tikz}
\usetikzlibrary{
  positioning, calc, arrows.meta, fit, backgrounds, shapes.geometric,
  shapes.misc, decorations.pathreplacing, decorations.pathmorphing,
  shadows.blur, patterns, math
}

% AUTO-GENERATED from palette.json by emit_palette.py. Do not edit by hand.
% Each figure sets \rltheme to light or dark before \input; the preamble loads this file.
\usepackage{etoolbox}
\providecommand{\rltheme}{light}%
\ifdefstring{\rltheme}{dark}{%
  \definecolor{rlInk}{HTML}{E6E8EB}%
  \definecolor{rlInkTwo}{HTML}{B7BDC7}%
  \definecolor{rlMuted}{HTML}{9AA0AA}%
  \definecolor{rlLine}{HTML}{3A3D44}%
  \definecolor{rlSurface}{HTML}{1D1E22}%
  \definecolor{rlSurfaceTwo}{HTML}{26272C}%
  \definecolor{rlBaseline}{HTML}{55565E}%
  \definecolor{rlTierObs}{HTML}{2DD4BF}%
  \definecolor{rlTierCausal}{HTML}{F59E0B}%
  \definecolor{rlTierVuln}{HTML}{FB7185}%
  \definecolor{rlTrustExp}{HTML}{64748B}%
  \definecolor{rlTrustCal}{HTML}{60A5FA}%
  \definecolor{rlTrustReg}{HTML}{A78BFA}%
  \definecolor{rlTrustAdj}{HTML}{22C55E}%
  \definecolor{rlSignal}{HTML}{E879F9}%
  \definecolor{rlGauge}{HTML}{22D3EE}%
  \definecolor{rlChosen}{HTML}{5598E7}%
  \definecolor{rlRejected}{HTML}{F0716F}%
  \definecolor{rlSeriesB}{HTML}{F98A5C}%
  \definecolor{rlAccent}{HTML}{FB8C5A}%
  \definecolor{rlSeqLo}{HTML}{0D366B}%
  \definecolor{rlSeqHi}{HTML}{CDE2FB}%
  \definecolor{rlDivMid}{HTML}{2A2B31}%
}{%
  \definecolor{rlInk}{HTML}{16181D}%
  \definecolor{rlInkTwo}{HTML}{4B5563}%
  \definecolor{rlMuted}{HTML}{6B7280}%
  \definecolor{rlLine}{HTML}{D7DAE0}%
  \definecolor{rlSurface}{HTML}{FFFFFF}%
  \definecolor{rlSurfaceTwo}{HTML}{F4F5F7}%
  \definecolor{rlBaseline}{HTML}{C3C2B7}%
  \definecolor{rlTierObs}{HTML}{0D9488}%
  \definecolor{rlTierCausal}{HTML}{B45309}%
  \definecolor{rlTierVuln}{HTML}{BE123C}%
  \definecolor{rlTrustExp}{HTML}{94A3B8}%
  \definecolor{rlTrustCal}{HTML}{3B82F6}%
  \definecolor{rlTrustReg}{HTML}{7C3AED}%
  \definecolor{rlTrustAdj}{HTML}{15803D}%
  \definecolor{rlSignal}{HTML}{A21CAF}%
  \definecolor{rlGauge}{HTML}{0891B2}%
  \definecolor{rlChosen}{HTML}{2A78D6}%
  \definecolor{rlRejected}{HTML}{E34948}%
  \definecolor{rlSeriesB}{HTML}{EB6834}%
  \definecolor{rlAccent}{HTML}{E85D2C}%
  \definecolor{rlSeqLo}{HTML}{CDE2FB}%
  \definecolor{rlSeqHi}{HTML}{0D366B}%
  \definecolor{rlDivMid}{HTML}{F0EFEC}%
}%
 % defines the rl* colors for the current \rltheme

% digit-name aliases, so both rlInkTwo and rlInk2 resolve
\colorlet{rlInk2}{rlInkTwo}
\colorlet{rlSurface2}{rlSurfaceTwo}

% ---- global defaults -------------------------------------------------------
\tikzset{
  >={Stealth[length=2.4mm, width=1.8mm]},
  font=\sffamily,
  every node/.append style={text=rlInk},
}

% ---- chips: the three-chip vocabulary (tier / gauge status / works-on) ------
% Outline chip reads correctly on both themes because the role color is a mid-tone.
\tikzset{
  rlchip/.style 2 args={
    draw=#1, text=#1, line width=0.5pt, rounded corners=2.5pt,
    inner xsep=5pt, inner ysep=2.6pt, font=\sffamily\bfseries\scriptsize,
    label/.expanded=#2,
  },
  rlchip/.default={rlMuted}{},
  % filled variant, surface-colored text so it contrasts in both themes
  rlchipfill/.style={
    draw=#1, fill=#1, text=rlSurface, rounded corners=2.5pt,
    inner xsep=5pt, inner ysep=2.6pt, font=\sffamily\bfseries\scriptsize,
  },
}
% simple outline chip taking one color (most common)
\newcommand{\rlchipnode}[3][rlMuted]{\node[draw=#1, text=#1, line width=0.6pt,
  rounded corners=2.5pt, inner xsep=5pt, inner ysep=2.6pt,
  font=\sffamily\bfseries\scriptsize] (#2) {#3};}

% ---- cards and panels ------------------------------------------------------
\tikzset{
  rlcard/.style={
    draw=rlLine, fill=rlSurface, line width=0.7pt, rounded corners=4pt,
    blur shadow={shadow blur steps=7, shadow xshift=0pt, shadow yshift=-1pt,
      shadow opacity=18, shadow blur radius=3pt},
  },
  rlpanel/.style={draw=rlLine, fill=rlSurfaceTwo, line width=0.6pt, rounded corners=3pt},
  rlfield/.style={font=\sffamily\scriptsize, text=rlMuted, align=left},
  rlvalue/.style={font=\sffamily\bfseries, text=rlInk},
  rllabel/.style={font=\sffamily\footnotesize, text=rlInk2},
  rlnote/.style={font=\sffamily\scriptsize\itshape, text=rlMuted},
}

% ---- lines, axes, vectors --------------------------------------------------
\tikzset{
  rlaxis/.style={draw=rlLine, line width=0.7pt, -{Stealth[length=1.8mm]}},
  rlbaseline/.style={draw=rlBaseline, line width=0.6pt, dash pattern=on 2pt off 2pt},
  rlflow/.style={draw=rlInk2, line width=0.9pt, -{Stealth[length=2.4mm]}},
  rlvec/.style={line width=1.3pt, -{Stealth[length=2.8mm, width=2.2mm]}},
  rldrop/.style={draw=rlMuted, line width=0.5pt, dash pattern=on 1.6pt off 1.6pt},
}

% ---- gate glyph: a narrow keyhole the evidence must pass to climb a rung ----
% \rlgate{name}{fill color}{label}  draws a small rounded gate node.
\newcommand{\rlgate}[3]{%
  \node[draw=#2, fill=rlSurface, line width=0.9pt, rounded corners=2pt,
        minimum width=6.5mm, minimum height=6.5mm, inner sep=0pt] (#1) {};
  \node[text=#2, font=\sffamily\bfseries\footnotesize] at (#1) {#3};%
}

% ---- answer-key glyph (calibration): a tiny key ----------------------------
\newcommand{\rlkey}[2][rlTrustAdj]{%
  \begin{scope}[shift={(#2)}, rotate=45]
    \draw[fill=#1, draw=#1] (0,0) circle (1.6mm);
    \fill[rlSurface] (0,0) circle (0.6mm);
    \draw[#1, line width=1.1pt] (0,-1.6mm) -- (0,-4.4mm);
    \draw[#1, line width=1.1pt] (0,-3.0mm) -- (1.1mm,-3.0mm);
    \draw[#1, line width=1.1pt] (0,-4.0mm) -- (0.9mm,-4.0mm);
  \end{scope}%
}

% ---- frozen-study seal: a small notched stamp ------------------------------
\newcommand{\rlseal}[2][rlTrustReg]{%
  \node[star, star points=12, star point ratio=0.9, draw=#1, fill=#1,
        text=rlSurface, minimum size=9mm, inner sep=0pt,
        font=\sffamily\bfseries\tiny] at (#2) {\#hash};%
}

\begin{document}
\begin{tikzpicture}[scale=1.0,
    % glyph vocabulary: neutral structure in ink/muted, family accent in rlSignal
    gline/.style={draw=rlInk2, line width=0.7pt},
    gmut/.style={draw=rlMuted, line width=0.6pt},
    gsig/.style={draw=rlSignal, line width=0.8pt},
    gbox/.style={draw=rlMuted, fill=rlSurfaceTwo, line width=0.6pt, rounded corners=1pt},
    % each grader taps the shared bus; the tap points toward the protocol
    rlattach/.style={draw=rlSignal, line width=0.9pt, -{Stealth[length=1.7mm]}},
    rlbus/.style={draw=rlSignal, line width=2pt, line cap=round,
      -{Stealth[length=3mm, width=2.6mm]}},
  ]
  \def\sy{1.5}   % the shared bus runs through the central gutter

  % subtle theme-matched backdrop (never a hardcoded white card)
  \fill[rlSurfaceTwo, rounded corners=8pt] (-1.75,-2.2) rectangle (14.8,4.35);

  % the destination: one protocol, drawn as the family-accented endpoint card
  \node[draw=rlSignal, fill=rlSurface, line width=1.15pt, rounded corners=5pt,
        minimum width=3.6cm, minimum height=4.75cm, inner sep=0pt] (panel) at (12.6,\sy) {};
  % header band
  \fill[rlSignal, rounded corners=3pt] (11.05,3.15) rectangle (14.15,3.72);
  \node[font=\sffamily\bfseries\small, text=rlSurface] at (12.6,3.44) {one protocol};
  \node[font=\sffamily\scriptsize, text=rlMuted] at (12.6,2.8) {every grader implements};
  % the three shared verbs
  \foreach \vy/\verb in {2.18/{score}, 1.45/{score over positions}, 0.72/{expose activations}}{
    \fill[rlSignal] (11.28,\vy) circle (0.05);
    \node[anchor=west, font=\sffamily\footnotesize, text=rlInk] at (11.46,\vy) {\verb};
  }
  \draw[rlLine, line width=0.6pt] (11.15,0.22) -- (14.05,0.22);
  \node[font=\sffamily\scriptsize, text=rlMuted] at (12.6,-0.14) {eight adapters};

  % the shared bus: eight taps collapse into one line into the protocol
  \draw[rlbus] (-0.55,\sy) -- (panel.west);

  % the eight grader cards, two rows of four
  \foreach \cx/\cy/\lbl in {%
      0/3.0/{classifier RM}, 2.9/3.0/{LLM judge}, 5.8/3.0/{process RM}, 8.7/3.0/{implicit (DPO)},
      0/0.0/{rubric head}, 2.9/0.0/{trajectory}, 5.8/0.0/{dense per-token}, 8.7/0.0/{ensemble}}{
    \node[rlcard, minimum width=2.5cm, minimum height=1.75cm, inner sep=0pt] at (\cx,\cy) {};
    \node[rllabel, text=rlInk2] at (\cx,\cy-0.55) {\lbl};
  }

  % taps: top row drops down onto the bus, bottom row rises up onto it
  \foreach \cx in {0,2.9,5.8,8.7}{
    \draw[rlattach] (\cx,2.125) -- (\cx,\sy+0.02);
    \draw[rlattach] (\cx,0.875) -- (\cx,\sy-0.02);
  }

  % ---- glyphs (upper half of each card; label sits below) ------------------
  % 1. classifier RM: a reward head on a stack of layers
  \begin{scope}[shift={(0,3.30)}]
    \foreach \yy in {-0.26,-0.12,0.02}{ \draw[gbox] (-0.34,\yy) rectangle (0.34,\yy+0.09); }
    \draw[rlSignal, line width=0.8pt, -{Stealth[length=1.3mm]}] (0,0.13) -- (0,0.185);
    \fill[rlSignal] (0,0.30) circle (0.10);
  \end{scope}

  % 2. LLM judge: a speech bubble delivering a verdict
  \begin{scope}[shift={(2.9,3.30)}]
    \draw[gline, fill=rlSurfaceTwo, rounded corners=2.5pt] (-0.42,-0.04) rectangle (0.42,0.34);
    \fill[rlSurfaceTwo] (-0.24,0.0) -- (-0.30,-0.22) -- (-0.10,0.0) -- cycle;
    \draw[gline] (-0.24,-0.005) -- (-0.30,-0.22) -- (-0.10,-0.005);
    \draw[rlSignal, line width=1.2pt, line cap=round] (-0.13,0.14) -- (-0.03,0.05) -- (0.16,0.25);
  \end{scope}

  % 3. process RM: rising steps, each with a per-step tick
  \begin{scope}[shift={(5.8,3.28)}]
    \foreach \i in {0,1,2}{
      \pgfmathsetmacro{\sx}{-0.38+\i*0.30}
      \pgfmathsetmacro{\syy}{-0.22+\i*0.13}
      \draw[gbox] (\sx,\syy) rectangle (\sx+0.20,\syy+0.15);
      \draw[rlSignal, line width=0.9pt, line cap=round]
        (\sx+0.04,\syy+0.22) -- (\sx+0.08,\syy+0.18) -- (\sx+0.15,\syy+0.30);
    }
  \end{scope}

  % 4. implicit (DPO): two policies, an implicit log-ratio readout
  \begin{scope}[shift={(8.7,3.30)}]
    \draw[gbox] (-0.48,0.06) rectangle (-0.06,0.30);
    \node[font=\tiny, text=rlInk2] at (-0.27,0.18) {$\pi_\theta$};
    \draw[gbox] (-0.48,-0.30) rectangle (-0.06,-0.06);
    \node[font=\tiny, text=rlMuted] at (-0.27,-0.18) {$\pi_{\mathrm{ref}}$};
    \draw[rlSignal, line width=0.9pt] (-0.02,0.18) -- (0.10,0.18) -- (0.10,-0.18) -- (-0.02,-0.18);
    \node[font=\tiny\bfseries, text=rlSignal, anchor=west] at (0.14,0) {log};
  \end{scope}

  % 5. rubric head: several criteria bars of differing length
  \begin{scope}[shift={(0,0.30)}]
    \draw[gmut] (-0.40,0.32) -- (-0.40,-0.30);
    \fill[rlSignal, fill opacity=0.90] (-0.40,0.19) rectangle (0.24,0.29);
    \fill[rlSignal, fill opacity=0.62] (-0.40,0.03) rectangle (0.02,0.13);
    \fill[rlSignal, fill opacity=0.78] (-0.40,-0.13) rectangle (0.16,-0.03);
    \fill[rlSignal, fill opacity=0.50] (-0.40,-0.29) rectangle (-0.06,-0.19);
  \end{scope}

  % 6. trajectory: a scored path through waypoints
  \begin{scope}[shift={(2.9,0.30)}]
    \draw[gline, line cap=round, -{Stealth[length=1.8mm]}]
      plot[smooth,tension=0.7] coordinates {(-0.45,-0.16)(-0.18,0.15)(0.08,-0.10)(0.30,0.13)(0.46,0.26)};
    \foreach \p in {(-0.45,-0.16),(-0.18,0.15),(0.08,-0.10),(0.30,0.13)}{
      \fill[rlSignal] \p circle (0.03); }
  \end{scope}

  % 7. dense per-token: a token strip with a bar over each token
  \begin{scope}[shift={(5.8,0.30)}]
    \foreach \i/\hh in {0/0.12,1/0.27,2/0.16,3/0.34,4/0.20,5/0.29}{
      \pgfmathsetmacro{\tx}{-0.45+\i*0.16}
      \draw[gbox] (\tx,-0.30) rectangle (\tx+0.12,-0.16);
      \fill[rlSignal] (\tx+0.02,-0.13) rectangle (\tx+0.10,-0.13+\hh);
    }
  \end{scope}

  % 8. ensemble: several graders stacked into one
  \begin{scope}[shift={(8.7,0.28)}]
    \draw[gmut, fill=rlSurfaceTwo, rounded corners=2pt] (-0.34,-0.28) rectangle (0.22,0.06);
    \draw[gline, fill=rlSurfaceTwo, rounded corners=2pt] (-0.26,-0.18) rectangle (0.30,0.16);
    \draw[gsig, fill=rlSurface, rounded corners=2pt] (-0.18,-0.08) rectangle (0.38,0.26);
    \node[font=\tiny\bfseries, text=rlSignal] at (0.10,0.09) {N};
  \end{scope}

  % the plain-words payoff (not a title; the mechanism behind the claim)
  \node[rlnote, text=rlMuted, anchor=north west, align=left] at (-1.45,-1.05)
    {Different objects, one interface. An instrument binds to the interface, not the grader,\\so the same measurement runs on all eight without a line of change.};
\end{tikzpicture}
\end{document}
